% ============================= Probes =============================
\section{How Much Orientation Is in the Frozen Vision Representation?}
\label{sec:probes}

The annotation analysis shows failures are not ``the model cannot see the
pose.'' But inspection is not measurement. We measure directly how much
object-intrinsic orientation is \emph{decodable} from the frozen Qwen2-VL-7B
vision features, with no language model in the loop.

\subsection{Protocol}
From the base model's vision tower we extract per-patch ViT embeddings (1280-d)
and post-merger features (3584-d --- the exact input the LLM receives), pooled
per image. Tasks: T1 facing vs.\ facing-away ($n{=}103$ test, majority 64.2\%),
T2 parallel vs.\ perpendicular ($n{=}34$, majority 64.7\%), T3 4-way
($n{=}137$, majority 49.6\%). Readouts: linear logistic regression, MLP
(256-hidden), and a \emph{patch-vote} probe (per-patch linear + majority vote,
preserving spatial structure). Training uses the audited-clean VSR train split
(no test examples); results are 5-fold CV and test accuracy with Wilson CIs.

\subsection{Results: orientation is weakly decodable}
Table~\ref{tab:probe} summarizes the test numbers (full cell-level table in
App.~C). On T1, every readout sits at or near the 64.2\% majority
baseline; balanced accuracy ranges 49--65\%. On T2, only the patch-vote probe
finds signal ($73.5\%$, $+8.8$ over majority) --- axis-alignment geometry is
partially readable from \emph{local} structure, while pooled readouts stay near
chance. On T3, everything is at chance. MLPs never beat linear probes: there is
no evidence of nonlinearly-accessible orientation information.

\begin{table}[t]
\centering
\footnotesize
\setlength{\tabcolsep}{4pt}
\begin{threeparttable}
\caption{\textbf{Frozen-feature probes on the 7B vision tower} (test accuracy,
\%, majority baseline in parentheses). Orientation is not robustly decodable;
only the patch-vote probe on the axis-alignment task (T2) shows real signal.}
\label{tab:probe}
\begin{tabular}{lccc}
\toprule
Readout & T1 f./away & T2 par./perp & T3 4-way \\
 & (maj.\ 64.2) & (maj.\ 64.7) & (maj.\ 49.6) \\
\midrule
ViT lin.\ (pooled)   & 65.0 & 61.8 & 52.6 \\
ViT MLP              & 55.3 & 67.6 & 52.6 \\
Merger lin.          & 68.9 & 64.7 & 53.3 \\
Merger MLP           & 52.4 & 64.7 & 47.4 \\
ViT patch-vote       & 62.1 & \textbf{73.5} & 51.8 \\
Merger patch-vote    & 60.2 & 67.6 & 52.6 \\
\bottomrule
\end{tabular}
\end{threeparttable}
\end{table}

\subsection{Object conditioning does not recover it}
A global pooled readout may be too weak: VSR orientation is
\emph{object-pair-relative} (``the cat is facing the dog''), and the probe does
not know which regions to compare. We therefore repeat the analysis with
\emph{object-grounded} readouts: subject/reference boxes from frozen Qwen2-VL
grounding (94.6\% of 647 orientation images have both boxes), region features
mean-pooled inside each box ($[\text{subj}, \text{ref}, \text{subj}-\text{ref},
\text{subj}\cdot\text{ref}]$), optionally concatenated with 12-d normalized box
geometry, fed to the same linear/MLP probes. Result (App.~C,
full table): object-conditioned \emph{visual} features do not jump --- T1 CV
stays 56--60\% $\approx$ majority, T2 and T3 at chance. The single above-chance
cell is the \emph{geometry-only} merger MLP ($71.7\%$ test), a box-statistics
signal from the grounding model's own boxes, whose CV (61.1\%) stays below
majority --- not a robust vision-feature signal.

\textbf{Interpretation.} The frozen visual representation encodes
object-intrinsic direction only weakly, and the absence is not an artifact of
global pooling. Notably, the generative 7B model exceeds every probe on the hard
\emph{facing} class ($73$--$75\%$), so useful signal can still emerge through
multimodal interaction or language priors --- which is exactly the component the
probes strip away. The bottleneck is therefore both representational (weak
decodability) and behavioral (what the generative interface does with it), which
motivates the interventions of the next section.
